\chapter{Testbed and Selectors}
\label{ca:chap:testbed}

\section{Experimental Setup}
\label{ca:sec:setup}

\paragraph{Domains.} We use three domains whose answers a program can check: three-number
Countdown puzzles \autocite{patino_gold_2025}, MATH problems at difficulty levels 3 to 5
\autocite{hendrycks_math_2021}, and Python function-writing problems from KodCode
\autocite{xu_kodcode_2025}, kept only if their reference solution passes its tests in our
sandbox. Each domain has a 2,000-prompt training pool, a 500-prompt selection set that
feeds the selector, and a 500-prompt test set used only for final results. Generated code
runs in a bubblewrap sandbox with no network and a read-only file system.

\paragraph{Student.} The student is Qwen2.5-1.5B-Instruct \autocite{qwen_qwen25_2024},
trained through a rank-16 LoRA adapter \autocite{hu_lora_2022}. Before training it scores
0.208 on Countdown, 0.474 on MATH and 0.396 on code on the selection set.

\paragraph{Providers.} Nine open models were screened on 200 prompts per domain and five
were kept (Table~\ref{ca:tab:pool}). They run locally but are reached only through a text
interface, so black-box access is enforced rather than assumed. The pool has a strong and
expensive model, a mid-priced generalist, a code specialist, a cheap model the student can
catch up with, and a cheap model from a second family, which separates any effect of
sharing the student's family from an effect of price. A sixth provider is a copy of Phi-4
at the same price whose answers are corrupted on 30\% of prompts.

\begin{table}[htbp]
\centering
\small
\begin{tabular}{@{}llrccc@{}}
\toprule
 & & & \multicolumn{3}{c}{\textbf{Ceiling $\mu_{k,d}$}} \\
\cmidrule(l){4-6}
\textbf{Provider} & \textbf{Size} & \textbf{Price $c_k$} & \textbf{Countdown} &
  \textbf{MATH} & \textbf{Code} \\
\midrule
Phi-4 \autocite{abdin_phi4_2024}                 & 14.7B & 2.92 & 0.874 & 0.804 & 0.686 \\
Qwen2.5-7B \autocite{qwen_qwen25_2024}           & 7.6B  & 1.00 & 0.810 & 0.750 & 0.580 \\
Qwen2.5-Coder-7B \autocite{hui_qwen25coder_2024} & 7.6B  & 0.98 & 0.620 & 0.678 & 0.612 \\
Qwen2.5-3B \autocite{qwen_qwen25_2024}           & 3.1B  & 0.50 & 0.538 & 0.690 & 0.462 \\
Phi-4-mini \autocite{microsoft_phi4mini_2025}    & 3.8B  & 0.84 & 0.616 & 0.618 & 0.558 \\
Phi-4, degraded                                  & 14.7B & 2.92 & 0.614 & 0.532 & 0.492 \\
\bottomrule
\end{tabular}
\caption{The strongest provider is also the most expensive. Price is GPU time per response
relative to Qwen2.5-7B; the ceiling is the accuracy of the provider's most common answer
over its four cached responses on the selection set.}
\label{ca:tab:pool}
\end{table}

\paragraph{Prices and cache.} A provider's price is the GPU time it needs to produce one
response, relative to Qwen2.5-7B, measured during screening and fixed before training.
Long-winded providers therefore cost more, as under per-token pricing. Each provider
answered every pool and selection prompt four times, once; runs read these cached answers,
which makes them reproducible and the measurements of Section~\ref{ca:sec:branching}
affordable.

\paragraph{Training and evaluation.} Every run has a budget of $B = 2{,}400$ price units
and a cap of 400 rounds. The selector chooses an arm once per block of 10 rounds, after
which the student is evaluated greedily on 256 selection prompts per domain. Final students
are evaluated on all 1,500 test problems. Every condition is run with three seeds, and
intervals are 95\% bootstrap intervals over seeds and test items.

\section{Selectors}
\label{ca:sec:selectors}

\paragraph{Headroom with learned efficiency (ours).} The selector scores each arm by how
far the student is below the provider's ceiling, scaled by a learned estimate of how much
accuracy that headroom buys, and divided by price:
\[
  \hat{v}(k, d) = \max\big(\hat{g}_{k,d},\, 0\big) \cdot
                  \frac{\max\big(\mu_{k,d} - \alpha_d,\ 0.05\big)}{c_k},
\]
where $\alpha_d$ is the student's current selection accuracy on domain $d$. The efficiency
$\hat{g}_{k,d}$ is the realised accuracy gain per round per unit of headroom, pooled per
domain, shrunk towards the domain's estimate and discounted by 0.95 per block, so that a
domain where no provider produces progress is learned to be unproductive for all of them at
once. With probability 0.1 the selector plays a random arm instead, so every arm keeps
being observed.

\paragraph{Cost-aware Thompson sampling.} A model-free comparison that keeps a discounted
Gaussian posterior over each arm's gain per round and plays the arm with the largest
sampled gain divided by price \autocite{xia_budgeted_ts_2015,raj_dts_2017}.

\paragraph{Reference policies.} Uniform allocation over all 18 arms; each real provider
alone, with domains in rotation; a control trained on ground-truth correctness with no
provider and no cost for 200 rounds, which shows what providers add beyond the labels; and
supervised fine-tuning (SFT) on provider responses bought with the same budget, which shows
what learning on-policy adds.

\section{Measuring Marginal Value}
\label{ca:sec:branching}

We measure $v(k, d; \theta)$ directly by branching. The uniform run of each seed saves
checkpoints at 25, 50 and 75\% of its budget, and the untrained student serves as the 0\%
checkpoint. From each checkpoint, a branch trains on a single arm for 20 rounds and is then
evaluated on the full selection set; marginal value is the change in mean selection
accuracy divided by the budget the branch spent. With 6 providers, 3 domains, 4 checkpoints
and 3 seeds this gives 216 branches. Branches from the same checkpoint and seed share the
adapter, the prompt order and the sampling seeds, so comparisons between providers are
paired.

\section{Design Decisions}
\label{ca:sec:decisions}

Table~\ref{ca:tab:decisions} lists the choices that most shape the results, the
alternatives we considered, and the evidence behind each choice.

\begin{table}[htbp]
\centering
\footnotesize
\renewcommand{\arraystretch}{1.1}
\begin{tabular}{@{}>{\raggedright\arraybackslash}p{2.7cm}
                >{\raggedright\arraybackslash}p{3.9cm}
                >{\raggedright\arraybackslash}p{8.3cm}@{}}
\toprule
\textbf{Decision} & \textbf{Alternatives} & \textbf{Reason for the choice} \\
\midrule
Black-box agreement reward & Logit-based on-policy distillation across tokenizers &
  Cross-tokenizer distillation recovered 0.14 of the teacher gap against 0.66 with a shared
  tokenizer, and its hybrid variant collapsed (Takeaway~\ref{ca:tk:blackbox}). \\
Normalised vote share & Most-common-answer indicator; raw vote share; rubric with a verifier &
  Carried as much learning signal as ground truth on MATH and code in screening, kept signal
  on Countdown where providers often give invalid answers, and needs no extra model. \\
Ceiling as provider quality & Single-response accuracy; agreement with other providers &
  It is where the reward points the student. Agreement with peers was label-free but
  favoured the Qwen family, which made up most of the pool. \\
Countdown (3 numbers), MATH 3--5, KodCode & GSM8K; MBPP; 4-number Countdown &
  GSM8K was saturated (all large providers scored 94--97\%), MBPP had too few problems, and
  the student solved 0.8\% of 4-number puzzles, leaving almost no gradient. \\
Price = GPU time per response & Parameter count; hosted API prices &
  Measured on our hardware and fixed before training, and captures verbosity: one candidate
  repeated itself until the length limit. \\
Corrupt whole prompts in the degraded provider & Corrupt each response independently &
  Independent corruption at 30\% lowered the ceiling by only 8\%, because the majority
  usually survives; per-prompt corruption lowers it by 30\%. \\
\bottomrule
\end{tabular}
\caption{Main design decisions and their evidence.}
\label{ca:tab:decisions}
\end{table}
